% Figura: ventanas y ventanas desplazadas del Swin Transformer. Adaptada de liu2021swin (Figura 2).
% Se usa en sec:marco_vit, en el párrafo del Swin.

\begin{figure}[H]
    \centering
    \begin{tikzpicture}[x=1mm, y=1mm]

        % Capa l: cuatro ventanas de 4x4 parches.
        \foreach \c/\x/\y in {figazul/0/16, fignaranja/16/16, figverde/0/0, figpurpura/16/0} {
            \fill[\c!22] (\x, \y) rectangle (\x + 16, \y + 16);
        }
        \draw[figtrazo!40] (0, 0) grid[step=4] (32, 32);
        \draw[figtrazo, very thick] (0, 0) rectangle (32, 32);
        \draw[figtrazo, very thick] (16, 0) -- (16, 32) (0, 16) -- (32, 16);
        \node[font=\footnotesize, text=figtrazo] at (16, -5) {capa $l$: ventanas regulares};

        % Capa l+1: ventanas desplazadas dos parches.
        \begin{scope}[xshift=50mm]
            \foreach \c/\x/\y/\w/\h in {
                figazul/0/24/8/8, fignaranja/8/24/16/8, figverde/24/24/8/8,
                figpurpura/0/8/8/16, figceleste/8/8/16/16, figbermellon/24/8/8/16,
                figazul/0/0/8/8, fignaranja/8/0/16/8, figverde/24/0/8/8} {
                \fill[\c!22] (\x, \y) rectangle (\x + \w, \y + \h);
            }
            \draw[figtrazo!40] (0, 0) grid[step=4] (32, 32);
            \draw[figtrazo, very thick] (0, 0) rectangle (32, 32);
            \draw[figtrazo, very thick] (8, 0) -- (8, 32) (24, 0) -- (24, 32) (0, 8) -- (32, 8) (0, 24) -- (32, 24);
            \node[font=\footnotesize, text=figtrazo] at (16, -5) {capa $l+1$: ventanas desplazadas};
        \end{scope}

        \draw[figflecha] (36, 16) -- (46, 16);

    \end{tikzpicture}
    \caption{Ventanas y ventanas desplazadas del Swin Transformer. Adaptada de \textcite{liu2021swin}.}
    \label{fig:marco_swin}
\end{figure}
